\chapter{Bladder Cancer}
\index{Bladder Cancer}

\section*{Definition and Overview}
Bladder cancer is a malignant growth that begins in the lining of the bladder, the hollow, muscular organ in the lower pelvis that stores urine before it leaves the body. The vast majority of bladder cancers arise from the urothelium, the specialised layer of cells that lines the inner surface of the bladder, ureters, and part of the kidneys. These are known as urothelial carcinomas (previously called transitional cell carcinomas). Less common types include squamous cell carcinoma, adenocarcinoma, and rare small-cell variants. A critical distinction governs behaviour and treatment: non-muscle-invasive bladder cancer (NMIBC) remains confined to the lining and has an excellent outlook with local treatment, whereas muscle-invasive bladder cancer (MIBC) has grown into the detrusor muscle wall and carries a significant risk of spread (metastasis) to lymph nodes and distant organs.

\section*{Epidemiology}
Bladder cancer is among the ten most common cancers worldwide and is significantly more common in men, who are affected roughly three to four times as often as women. It is predominantly a disease of older age, with most cases diagnosed after the sixth decade; the median age at diagnosis is around 70 years. Incidence varies widely between countries, closely tracking historical smoking rates and exposure to industrial carcinogens. In high-income countries, roughly 75\% of new cases are non-muscle-invasive at diagnosis. Five-year survival is excellent for localised disease but falls sharply once the tumour has invaded muscle or spread, making early detection and surveillance priorities.

\section*{Aetiology and Pathophysiology}
The bladder lining is continuously bathed by urine, which concentrates and stores many of the chemicals filtered from the bloodstream. This explains why so many bladder-cancer risk factors are environmental and why prolonged exposure matters. Known risk factors include:

\begin{itemize}
    \item \textbf{Smoking:} the single most important risk factor, accounting for roughly half of cases; carcinogens in tobacco smoke are concentrated in urine and damage urothelial DNA
    \item \textbf{Occupational exposures:} aromatic amines and other chemicals used in dye, rubber, leather, textile, and paint industries
    \item \textbf{Chronic bladder irritation:} long-term indwelling catheters, recurrent infections, and chronic bladder stones, which predispose to squamous cell carcinoma
    \item \textbf{Pelvic radiation therapy} given for other cancers, especially gynaecological or prostate malignancy
    \item \textbf{Certain drugs:} cyclophosphamide and the diabetes medicine pioglitazone have been linked to increased risk
    \item \textbf{Genetic factors:} a family history of bladder cancer and, rarely, inherited syndromes such as Lynch syndrome
    \item \textbf{Arsenic:} chronic exposure through contaminated drinking water in some regions
\end{itemize}

In carcinogenesis, damaged urothelial cells acquire progressive mutations in oncogenes and tumour-suppressor genes. Most NMIBCs are low-grade, papillary growths that recur locally but rarely invade. High-grade tumours and carcinoma-in-situ carry a much higher risk of progressing to muscle invasion, after which lymphatic and haematogenous spread to the pelvic lymph nodes, bones, lungs, and liver becomes possible.

\section*{Clinical Features}
Painless visible blood in the urine (gross haematuria) is the classic presenting symptom and is the feature that most often leads to diagnosis. However, symptoms can be subtle or intermittent, and bleeding may only be detected on a urine dipstick test.

\begin{itemize}
    \item \textbf{Haematuria:} blood in the urine that may be visible or microscopic, often painless, and frequently comes and goes
    \item \textbf{Irritative urinary symptoms:} urgency, urinary frequency, and a burning sensation on urination that mimic a urinary tract infection
    \item \textbf{Pelvic or lower back pain}, sometimes extending to the flank
    \item \textbf{Recurrent urinary tract infections} that fail to resolve or keep returning
    \item \textbf{Advanced disease features:} unintentional weight loss, fatigue, bone pain, swelling of the legs, or a palpable lower-abdominal mass
\end{itemize}

\section*{Differential Diagnosis}
The symptoms of bladder cancer overlap with many benign conditions, and blood in the urine must never be assumed to be trivial, particularly in older people and smokers.

\begin{itemize}
    \item Urinary tract infection or acute cystitis
    \item Bladder stones or ureteric (kidney) stones
    \item Benign prostatic enlargement in men
    \item Prostate cancer or kidney (renal) cancer
    \item Glomerulonephritis and other intrinsic kidney disease
    \item Bleeding due to anticoagulant medication
    \item Schistosomiasis (bilharzia) in endemic regions
\end{itemize}

\section*{When to Seek Medical Care}
Any episode of blood in the urine -- even a single episode that clears within a day -- warrants prompt medical assessment. Similarly, seek review for persistent urinary symptoms or recurrent infections that do not settle despite treatment.

\textbf{Contact your local emergency services for:}
\begin{itemize}
    \item Inability to pass urine (acute urinary retention) with a painful, distended bladder
    \item Heavy visible bleeding with clots that interferes with urination
    \item Sudden, severe abdominal or flank pain
    \item High fever with shaking chills, which may indicate urosepsis, especially after procedures
\end{itemize}

\section*{Diagnosis and Investigations}
The diagnosis and staging of bladder cancer rely on a combination of urine tests, endoscopic inspection, and imaging, coordinated through a rapid-access referral pathway for visible haematuria.

\begin{itemize}
    \item \textbf{Urine testing:} dipstick analysis for blood, microscopy, and urine cytology to detect malignant cells shed from the urothelium
    \item \textbf{Cystoscopy:} the definitive investigation; a thin fibre-optic telescope passed through the urethra directly visualises the bladder lining
    \item \textbf{Biopsy and transurethral resection:} tissue is removed during cystoscopy, allowing confirmation of the diagnosis, the grade, and whether muscle is invaded
    \item \textbf{CT urogram:} contrast-enhanced imaging of the kidneys, ureters, and bladder, which identifies tumours and assesses the upper urinary tract
    \item \textbf{Ultrasound:} useful as a first-line test for kidney assessment and to screen for a bladder mass
    \item \textbf{Staging scans:} CT of the chest, abdomen, and pelvis, and sometimes PET imaging or MRI, when muscle-invasive disease is suspected
    \item \textbf{Risk stratification:} tumour stage (T), grade, and the presence of carcinoma-in-situ determine the surveillance and treatment plan
\end{itemize}

\section*{Management}
Treatment is planned by a multidisciplinary team and depends chiefly on whether the tumour has invaded the muscle wall, together with the tumour grade, the person's fitness, and their preferences.

\begin{itemize}
    \item \textbf{Non-muscle-invasive disease:} transurethral resection of the bladder tumour, followed by a single post-operative dose of intravesical chemotherapy; patients at higher risk receive repeated instillations of the immunotherapy agent BCG
    \item \textbf{Muscle-invasive disease:} radical cystectomy (removal of the bladder) with urinary diversion, or radical radiotherapy often combined with chemotherapy (trimodality therapy) as a bladder-sparing alternative
    \item \textbf{Neoadjuvant and adjuvant chemotherapy:} given before or after surgery in suitable patients to reduce the risk of recurrence
    \item \textbf{Advanced and metastatic disease:} systemic chemotherapy, immune checkpoint inhibitors, and targeted therapies that may control disease for prolonged periods
    \item \textbf{Urinary diversion:} after cystectomy, the urine is redirected through an ileal conduit, a continent pouch, or a bladder reconstruction using bowel
    \item \textbf{Surveillance:} regular cystoscopy and imaging, because recurrence is common even after successful initial treatment
    \item \textbf{Supportive care:} symptom control, stoma care, pelvic floor rehabilitation, and psychological support at every stage
\end{itemize}

\section*{Complications}
\begin{itemize}
    \item Local recurrence after transurethral resection, requiring repeated procedures
    \item Progression of non-muscle-invasive disease to muscle invasion and metastasis
    \item Spread to lymph nodes, bones, lungs, and liver in advanced disease
    \item Obstruction of the ureters, leading to hydronephrosis and kidney failure
    \item Bleeding, perforation, or infection after cystoscopic procedures
    \item Bladder irritation and haematuria from intravesical BCG or chemotherapy
    \item Long-term consequences of cystectomy, including altered body image, sexual dysfunction, and complications of the urinary diversion
\end{itemize}

\section*{Prognosis and Outcomes}
Outcomes are strongly determined by stage and grade at diagnosis. Non-muscle-invasive bladder cancer is usually controlled for many years, although it frequently recurs and therefore demands lifelong surveillance; progression to invasive disease is uncommon in low-grade tumours. Muscle-invasive cancer requires aggressive treatment, and cure is achievable when disease remains localised. Once the cancer has metastasised, cure is rarely possible, but modern systemic therapies can control disease for prolonged periods. Stopping smoking, attending scheduled follow-up, and prompt reporting of any new haematuria all improve the chances of a favourable outcome.

\section*{Prevention and Self-Care}
\begin{itemize}
    \item Do not smoke, and seek support to stop if you do; risk declines over time after cessation
    \item Follow workplace safety rules and use protective equipment in industries with chemical exposures
    \item Maintain a high fluid intake to keep the urine dilute and flush the bladder regularly
    \item Report any blood in the urine promptly, even if it clears spontaneously
    \item Seek advice for persistent urinary symptoms or repeated infections rather than ignoring them
    \item Attend all follow-up appointments and surveillance cystoscopies without fail
\end{itemize}

\section*{Sources and Further Reading}
The following sources provide reliable, up-to-date information for patients and students of medicine.

\begin{itemize}
    \item NHS. \textit{Overview of bladder cancer, its symptoms, and treatment.} https://www.nhs.uk/conditions/bladder-cancer/
    \item Mayo Clinic. \textit{Bladder cancer -- symptoms, causes, and treatment.} https://www.mayoclinic.org/diseases-conditions/bladder-cancer/
    \item MSD Manual. \textit{Bladder cancer -- professional edition.} https://www.msdmanuals.com/
    \item MedlinePlus. \textit{Bladder cancer -- patient education and resources.} https://medlineplus.gov/bladdercancer.html
    \item Centers for Disease Control and Prevention. \textit{Bladder cancer information and statistics.} https://www.cdc.gov/cancer/bladder/
\end{itemize}
